%% This is file `jcomp-template.tex',
%% 
%% Copyright 2017 Elsevier Ltd
%% 
%% This file is part of the 'Elsarticle Bundle'.
%% ---------------------------------------------
%% 
%% It may be distributed under the conditions of the LaTeX Project Public
%% License, either version 1.2 of this license or (at your option) any
%% later version.  The latest version of this license is in
%%    http://www.latex-project.org/lppl.txt
%% and version 1.2 or later is part of all distributions of LaTeX
%% version 1999/12/01 or later.
%% 
%% The list of all files belonging to the 'Elsarticle Bundle' is
%% given in the file `manifest.txt'.
%% 
%% Template article for Elsevier's document class `elsarticle'
%% with harvard style bibliographic references
%%
%% $Id: jcomp-template.tex 100 2017-07-14 13:15:12Z rishi $
%%
%% Use the option review to obtain double line spacing
%\documentclass[times,review,preprint,authoryear]{elsarticle}

%% Use the options `twocolumn,final' to obtain the final layout
%% Use longtitle option to break abstract to multiple pages if overfull.
%% For Review pdf (With double line spacing)
%\documentclass[times,twocolumn,review]{elsarticle}
%% For abstracts longer than one page.
%\documentclass[times,twocolumn,review,longtitle]{elsarticle}
%% For Review pdf without preprint line
%\documentclass[times,twocolumn,review,nopreprintline]{elsarticle}
%% Final pdf
\documentclass[times,final]{elsarticle}
%%
%\documentclass[times,twocolumn,final,longtitle]{elsarticle}
%%


%% Stylefile to load JCOMP template
\usepackage{jcomp}
\usepackage{framed,multirow}

%% The amssymb package provides various useful mathematical symbols
\usepackage{amssymb}
\usepackage{latexsym}

% Following three lines are needed for this document.
% If you are not loading colors or url, then these are
% not required.
\usepackage{url}
\usepackage{xcolor}
\definecolor{newcolor}{rgb}{.8,.349,.1}

\usepackage{bm}% bold math

\usepackage{todonotes}
\usepackage[acronym]{glossaries}
% \input{acronyms}

\newtheorem{lemma}{Lemma}
\newtheorem{corollary}{Corollary}

\DeclareMathOperator*{\argmax}{arg\,max}
\DeclareMathOperator*{\argmin}{arg\,min}

\journal{Building and Environment}

\begin{document}
	
	\verso{Alexander Hagg \textit{et al.}}
	
\begin{frontmatter}

\title{OpenSKIZZE: A Multi-Layered AI Framework for Co-Creative, Climate-Adaptive Urban Design}
% OpenSKIZZE: An Interactive Open-Source Tool for Exploring Climate-Resilient Urban Planning Strategies
% OpenSKIZZE: an AI assistant to Discover Climate Adapted Urban Planning Solutions
% Optimization-Driven Discovery of Context-Aware Urban Morphology
%\title{Data-Driven Discovery of Context-Aware Urban Morphology for Nighttime Cooling}

\author[1]{Alexander \snm{Hagg}\corref{cor1}}
\cortext[cor1]{Corresponding author: 
	Grantham-Allee 20; 53757 Sankt Augustin; Germany;}
\ead{alexander.hagg@h-brs.de}
\author[1,2]{Dirk \snm{Reith}}

\address[1]{Institute of Technology, Resource and Energy-Efficient Engineering, Bonn-Rhein-Sieg University of Applied Sciences, Sankt Augustin, Germany}
\address[2]{Fraunhofer Institute for Algorithms and Scientific Computing (SCAI), Sankt Augustin, Germany}

%\received{1 May 2013}
%\finalform{10 May 2013}
%\accepted{13 May 2013}
%\availableonline{15 May 2013}
%\communicated{S. Sarkar}


\begin{abstract}
Urban areas face significant challenges from climate change, particularly from increased heat stress. While advanced climate models exist, a gap persists in translating their insights into actionable, early-stage urban design processes. This paper introduces OpenSKIZZE, a novel, open-source framework that facilitates a co-creative workflow between designers and a multi-layered AI assistant. The foundation of OpenSKIZZE is a Quality-Diversity (QD) algorithm that comprehensively maps the solution space of urban designs, evaluating performance using the KLAM\_21 climate simulation tool and diverse morphological metrics. This foundational archive is then used to train two specialized deep learning models for interactive use: (1) a U-NET that provides instantaneous, high-fidelity predictions of airflow patterns for any user-drawn design, and (2) a conditional generative adversarial network (CGAN) that synthesizes novel, high-performing design options based on user-specified constraints. Through a web-based interface, users can explore the pre-computed archive, receive real-time predictive feedback on their own sketches, and generate new concepts on demand. We demonstrate and validate the framework through a real-world case study in Bonn, Germany, aimed at improving a vital cold-air corridor. Our validation confirms a strong correlation between the U-NET's predictions and ground-truth KLAM\_21 simulations. The results show that this multi-layered AI approach successfully bridges the praxis gap, enabling stakeholders to rapidly explore, understand, and create more resilient and climate-adaptive urban environments.

%This paper introduces an application for data-driven exploration of urban morphology tailored for nighttime cooling. Utilizing 3D voxel encodings, we investigate urban form impact on nighttime airflow via the KLAM\_21 model. Evolutionary quality diversity optimization generates diverse urban designs, representing qualitative variations across different environments. Location specificity is addressed by simulating airflow in three(?) distinct virtual environments. Gaussian process regression models enhance computational efficiency. The resulting design space, defined by general morphological features as niching dimensions, identifies high-performing designs that maximize wind speed for natural ventilation. This approach reveals both expected and unexpected context-aware design solutions. The open-source software and trained models are publicly available, facilitating further research and application for climate-resilient urban development.
\end{abstract}

\begin{keyword}
	%% MSC codes here, in the form: \MSC code \sep code
	%% or \MSC[2008] code \sep code (2000 is the default)
	%\MSC 41A05\sep 41A10\sep 65D05\sep 65D17
	%\MSC 68T20\sep 65D17\sep 68W50\sep76-10
	%% Keywords
	%\KWD Keyword1\sep Keyword2\sep Keyword3
	\KWD multi-solution optimization\sep machine learning\sep efficient domain analysis\sep problem diversity\sep generative methods\sep fluid dynamics
\end{keyword}

\end{frontmatter}

\section{Introduction}

Urbanization and climate change pose significant challenges to sustainable city development, particularly concerning heat accumulation and air quality. Natural ventilation is crucial for mitigating the urban heat island effect, enhancing air circulation, and improving thermal comfort. Traditional urban planning often relies on established guidelines and heuristic-based designs, which may limit the exploration of innovative solutions. However, the relationship between general and localized guidelines warrants closer examination. While general guidelines provide broad principles applicable across diverse contexts, significant discrepancies often arise when these principles are applied to specific locations with unique climatic, cultural, or spatial constraints. For instance, urban planning studies have documented how localized wind patterns in coastal cities or arid regions may require significantly different building orientations compared to general guidelines\todo{SOURCE}. Additionally, cultural preferences for open spaces or specific architectural styles can further challenge the applicability of broad design principles. 
On top of this, regional climate analysis, leading to concrete guidelines for specific districts or regions are based on urban climate simulations that are not able to resolve the specific air flow around buildings~\cite{Kruse2018lanuv}. 
So to actually discover and understand specific guidelines for concrete local urban planning in an automated manner is necessary, due to the large variety of cities and landscapes our planet and society has created.
This highlights the need for adaptable frameworks that respect both the universality of general guidelines and the specificity of localized contexts. 

Understanding how general frameworks translate—or fail to translate—into actionable solutions for localized contexts is essential to advancing urban planning practices. This insight can deepen critiques of traditional methods and highlight the importance of context-sensitive approaches. Optimization methods typically focus on single-objective or multi-objective approaches but often lack the capability to explore diverse and unconventional configurations comprehensively.

The main objective of this paper is to introduce a semi-automated framework that is able to discover localized design guidelines for climate resilient urban planning. We introduce a novel approach that combines Quality Diversity (QD) optimization and Reynolds-Averaged Navier-Stokes (RANS) simulations in OpenFOAM to autonomously discover implicit urban design guidelines. Unlike traditional methods, which are constrained by prior knowledge or subjective bias, our framework explores a vast design space defined by generalized morphological and flow features that are both relevant for urban planners and directly derivable from RANS simulations. We derive a machine learning model that is able to capture the patterns in the optimization landscape, allowing reuse of the localized implicit guidelines during the planning process.

While the framework minimizes human intervention, it is important to acknowledge the implicit biases embedded in the selection of features and objectives. These biases are informed by the necessity to balance generalization with practical applicability, ensuring that the discovered guidelines are meaningful and actionable for urban planning contexts. This proof-of-concept study highlights the potential of computational optimization tools to uncover both expected and unexpected insights, with the ultimate goal of complementing established urban planning practices.

\section{Literature}
\subsection{General Guidelines for Climate-Resilient Urban Planning}

General urban planning guidelines provide universal principles aimed at fostering sustainable and climate-resilient cities. For instance, the Architectural Institute of Japan (AIJ) provides standards to evaluate pedestrian wind comfort, while the Hong Kong Air Ventilation Assessment (AVA) Guidelines emphasize maximizing wind flow through urban areas to reduce heat accumulation. Similarly, the European EN 16798-1 standards and ASHRAE guidelines address ventilation and thermal comfort parameters that are globally applicable. These guidelines are valuable for establishing baseline design criteria, but they often assume homogeneity in environmental and cultural contexts.


However, localized factors such as unique wind patterns, topography, and cultural preferences can diverge significantly from these general principles. For example, in hot-arid climates, urban density and wind direction critically influence air velocity patterns, requiring tailored approaches that general guidelines do not explicitly account for. 

``However, there are hardly any studies on applied urban climatology in areas with a Mediterranean type of climate''~\cite{alcoforado2009application}, points to the problem that urban design guidelines for areas with specific properties often do not exist. 

The climate analysis report of the city of Bonn, Germany states that ``Diese idealtypischen Muster werden im gesamtstädtischen Kontext durch komplexe lokal-klimatische Effekte (nachbarschaftliche Wirkungen, horizontale und vertikale Strömungsprozesse) überprägt und können sich somit im konkreten räumlichen Fall auch (ganz) anders darstellen. Dennoch sind die skizzierten Phänomene grundlegend für das Verständnis des Modells und seiner Ergebnisse.'' idealization can help to understand urban climate effects in general, but lack specificity to the large variety of urban patterns ~\cite{berchtold2023MUTABOR}

Examples from studies such as Oke et al. (Urban Climates, Cambridge University Press, 2017???) and Santamouris' work on urban heat mitigation (Environmental Design of Urban Buildings, Earthscan Publications, 2006???) demonstrate how regional and local adaptations are necessary to make general principles actionable.

``Hierbei gilt es mit der Gebäudestruktur auf nächt-
liche Kaltluftströmungen, auf die Durchlüftung am Tag und auf das Zusammenspiel mit Grünräu-
men zu reagieren. Gebäudekubatur, insbesondere ein Bauen in die Höhe statt in die Fläche,
Gebäudestellung, -ausrichtung und -körnung können das Lokalklima beeinflussen. Da die In-
tegration dieser Maßnahmen in die Potenzialermittlung nicht automatisiert erfolgen kann und das
Forschungsprojekt einen Fokus auf den Gebäudebestand legt, werden diese Maßnahmen nur in
exemplarischen Teilräumen in die Potenzialermittlung integriert. Für drei Testgebiete wird die
Gebäudestruktur durch einen händischen Entwurf klimatisch optimiert [...]''~\cite{berchtold2023MUTABOR}

``Abhängig von den Bebauungsstrukturen kann diese bis zu 700 Meter (m) in die Siedlungsräume hineinwirken.'' in this climate analysis for the region of Northrhine-Westphalia in Germany, the raster size of the simulation is set to 100 m. x 100 m., which only allows a parameterized representation of the effects of smaller entities like buildings and trees~\cite{Kruse2018lanuv}.

``Lokal kann es zwar durch Düseneffekte auch zu einer Beschleunigung des Windes kommen, in der
summarischen Betrachtung über eine Gebäudebeinhaltende Rasterzelle überwiegt indes die Verlangsamung der Strömung. Gleichzeitig wird durch die Vielzahl der unterschiedlichen Strömungshindernisse die Turbulenz verstärkt. Darüber hinaus wird auch die Temperaturverteilung in starkem Maße modifiziert, da die in die bodennahe Atmosphäre ragenden Baukörper bis zur mittleren Bauhöhe in einem Wärmeaustausch mit der
Umgebung stehen. Die Temperatur wird durch die gebäudespezifischen Parameter wie Gebäudehöhe, Überbauungsgrad oder anthropogen Abwärme bestimmt und damit das Temperaturfeld der bodennahen Atmosphäre bis in die mittlere Höhe der Bebauung
modifiziert''~\cite{Kruse2018lanuv}



\subsection{Data-Driven Discovery in General Problem Solving}

Data-driven discovery methods have gained traction across various disciplines for their ability to uncover patterns and optimize solutions without relying on human intuition or preconceived notions. For instance, in materials science, data-driven methods have been used to optimize alloy compositions and predict material properties, as demonstrated in studies such as Fischer et al.'s work on computational material science~\cite{fischer2006predicting}. Similarly, computational biology employs optimization algorithms to predict protein-ligand interactions and guide drug discovery, exemplified by studies like Jumper et al. on AlphaFold's groundbreaking protein structure prediction~\cite{jumper2021highly}. These methods demonstrate the capacity to autonomously explore large and complex solution spaces, providing insights that traditional approaches might overlook.



\subsection{Data-Driven Discovery of Urban Planning Guidelines}

In the context of urban planning, data-driven methods are increasingly being employed to address challenges like traffic optimization, pollution control, and energy efficiency. For example, machine learning models trained on urban datasets can predict environmental impacts or identify effective mitigation strategies. Studies such as Chen et al. (Journal of Cleaner Production, 2020???) highlight how integrating data analytics with urban planning enhances air quality management. 

% Combining computational fluid dynamics and neural networks to characterize microclimate extremes: Learning the complex interactions between meso-climate and urban morphology
Combining fluid dynamics and neural networks to characterize interactions between the meso-climate and urban morphology~\cite{javanroodi2022combining}
The authors use nine morphological parameters but a fixed grid of buildings to predict wind speed and air temperature.


Furthermore, the application of computational fluid dynamics (CFD) with optimization algorithms has shown promise in exploring urban layouts that enhance natural ventilation, as demonstrated in Allegrini et al.~\cite{allegrini2015coupled}.

However, the novelty of our approach lies in the integration of Quality Diversity (QD) optimization with Reynolds-Averaged Navier-Stokes (RANS) simulations. This combination enables the autonomous discovery of localized urban design guidelines, bridging the gap between general principles and specific environmental or cultural needs. By combining evolutionary algorithms with RANS-derived flow features, the framework expands the scope of possible urban configurations while aligning with established principles. For example, configurations discovered using this approach can be cross-referenced with standards like the AIJ guidelines or Hong Kong AVA requirements to ensure applicability while fostering innovative solutions.

Our framework, as a proof of concept, demonstrates the potential to align discovered configurations with established standards, while also offering unexpected and innovative insights that traditional methods may overlook. The emphasis on exploring the full design space provides a robust mechanism for tailoring general principles to localized contexts.

\section{Methodology}
\todo[inline]{adequate scale of analysis}

\subsection{Design Generation}

\paragraph{Quality Diversity Optimization}

We employ an evolutionary QD optimization algorithm to generate diverse
urban layouts. These designs vary in terms of building shapes,
positions, dimensions, and orientations. The QD algorithm operates
within a design space defined by morphological features such as Building
Coverage Ratio (BCR), Floor Area Ratio, average building height,
building height variability, wind porosity, Number of Buildings, Average
Spacing Between Buildings, Open Space Ratio, and Urban Porosity.


\paragraph{Efficient Encodings}

We employ parameterizations that allow for efficient exploration of the
design space, encoding building configurations with fundamental
parameters like positions, dimensions, and orientations. This approach
minimizes the dimensionality of the search space while maximizing the
diversity of achievable designs.


\paragraph{Unbiased Features}


Morphological features such as Building Coverage Ratio (BCR), average
building height, and building height variability are used as niching
dimensions, defined over the entire urban plan. These features are
fundamental and directly measurable, avoiding subjective interpretation
or bias, and represent the entire design. They are universally applicable regardless of urban planning style.

\begin{itemize}
	\item Plot area ration: total floor area of buildings divided by total site area
	\item Site coverage index: total footprint divided by total site area
	\item Volume area ratio: total volume divided by site area
	\item Frontal area density: frontal area over the total horizontal surface area, depending on approaching wind direction
	\item Height to width ratio: aspect ratio of street canyon (???)
	\item Sky view factor: fraction of overlaying hemisphere occupied by the sky at each level
	\item Normalized height: of each point (w. r. t. maximum height)
	\item Effective length: \dots`
\end{itemize}~\cite{javanroodi2022combining}

Flow features like average wind speed, turbulence intensity, and
recirculation zones are calculated directly from simulation outputs.
These features inform the understanding of how different designs impact
airflow patterns.


\paragraph{Surrogate Models}

To accelerate the optimization process, we implement Gaussian process
regression (GPR) as surrogate models. These models are trained on data
generated during the evolutionary QD optimization process. The surrogate
models predict flow features such as average wind speed, turbulence
intensity, and recirculation zone areas, reducing the reliance on
computationally expensive RANS simulations.


\subsection{Simulation Setup}


\paragraph{OpenFOAM RANS Simulations}

RANS simulations in OpenFOAM form the foundation for evaluating flow
characteristics. We employ turbulence models like k-$\epsilon$ and k-$\omega$ SST, which
are well-suited for urban airflow simulations. The boundary conditions
include:

\begin{itemize}
\item
  Inflow wind profiles representative of realistic urban environments.
\item
  No-slip conditions on building surfaces.
\item
  Appropriate outlet conditions to ensure accurate flow predictions.
\end{itemize}


\subsection{Objective Function}

The primary objective is to maximize average wind speed at pedestrian
level (1.5 meters above ground) while maintaining a diverse set of
solutions. This ensures both improved natural ventilation and
exploration of varied urban configurations.


\subsection{Automation}

All aspects of the simulation process---from design generation to data
extraction and integration with the optimization algorithm---are
automated using custom scripting. This enables efficient exploration and
evaluation of the design space.


\section{Validation}

To validate the discovered guidelines, we compare our results against
established urban design standards:

\begin{enumerate}
\def\labelenumi{\arabic{enumi}.}
\item
  \textbf{Architectural Institute of Japan (AIJ):}

  \begin{itemize}
  \item
    Guidelines for evaluating wind effects on pedestrians.
  \end{itemize}
\item
  \textbf{Hong Kong Air Ventilation Assessment (AVA) Guidelines:}

  \begin{itemize}
  \item
    Focused on ventilation performance in dense urban environments.
  \end{itemize}
\item
  \textbf{European EN 16798-1 Standards:}

  \begin{itemize}
  \item
    Defines indoor environmental input parameters, including urban
    ventilation and thermal comfort.
  \end{itemize}
\item
  \textbf{ASHRAE Standards:}

  \begin{itemize}
  \item
    Comprehensive thermal environmental conditions for human occupancy.
  \end{itemize}
\item
  \textbf{Urban Climate and Design Resources:}

  \begin{itemize}
  \item
    Such as Oke et al.'s \emph{Urban Climates}, which provides a broader
    theoretical foundation for urban planning practices.
  \end{itemize}
\end{enumerate}

Quantitative comparisons involve overlaying simulated flow metrics
(e.g., wind speed, turbulence) with thresholds or recommendations
defined in these guidelines. This approach demonstrates how our
framework's outputs align with or extend existing urban planning
principles.


\section{Results}


\subsection{Optimization Outcomes}


\paragraph{Enhanced Ventilation Designs}

The framework successfully identifies urban configurations that
significantly improve natural ventilation efficiency at pedestrian
level.


\paragraph{Comprehensive Exploration}

By performing full domain analysis, we explore a wide range of design
possibilities, uncovering high-performing solutions across the solution
space.

\paragraph{Emergent Design Guidelines}


\subparagraph{Expected Findings}

The optimization process reveals known effective strategies, such as
optimal building spacing, alignment with wind direction, and appropriate
height variations, confirming the validity of the framework.


\subparagraph{Unexpected Insights}

Novel arrangements, including non-linear street layouts, asymmetrical
building distributions, and innovative building shapes, emerge as
effective strategies for enhancing ventilation, showcasing the
framework's ability to discover unconventional solutions.


\subsection{Validation Insights}

Comparison with established guidelines reveals areas of alignment and
divergence:


\paragraph{Alignment}

Many configurations adhere to principles such as wind corridor alignment
and optimal spacing for airflow.


\paragraph{Divergence}

Novel configurations suggest alternative strategies, such as leveraging
asymmetry to disrupt recirculation zones, which are not explicitly
addressed in existing guidelines.


\section{Discussion}

The QD optimization framework successfully identifies urban
configurations that enhance natural ventilation. The discovered
guidelines include both expected principles---such as optimal building
spacing and height variations---and unexpected insights, like the
effectiveness of non-linear street layouts and asymmetrical building
distributions.

\subsection{Impact on Urban Planning}

This framework demonstrates the potential of computational tools to
complement traditional urban planning by uncovering innovative solutions
through data-driven processes. By exploring a generalized feature space,
the approach ensures applicability across diverse urban contexts.


\subsection{Data-Driven Design}

The automated discovery of design guidelines underscores the potential
of computational tools in informing sustainable urban development,
providing empirical evidence to support decision-making.


\subsection{Unbiased Exploration}

By avoiding human bias and prior assumptions, the framework reveals
innovative solutions that might be overlooked in traditional planning,
fostering creativity and innovation in urban design.

\subsection{Methodological Contributions}


\paragraph{Framework for FDA in Urban Design}

We demonstrate how FDA principles can be applied to urban planning,
integrating optimization, simulation, and machine learning to understand
and navigate the full design space.


\paragraph{Efficiency in High-Dimensional Spaces}

The use of surrogate models and efficient encodings addresses the
challenges of computational cost in exploring large design spaces,
making the approach feasible for practical applications.

\subsection{Limitations}

As a proof of concept, this study relies on simplified assumptions, such
as steady wind conditions and uniform building materials. Future work
should incorporate additional environmental factors, such as pollutant
dispersion, thermal comfort, and dynamic wind conditions.

\paragraph{Computational Resources}

While surrogate modeling reduces computational demands, high-fidelity
simulations remain resource-intensive. Future work may focus on further
efficiency improvements, such as advanced surrogate modeling techniques
or parallel computing strategies.

\section{Future Work}

Further validation is needed through wind tunnel experiments or
real-world field studies. Additionally, integrating high-fidelity
simulations like LES and extending the framework to include
multidisciplinary objectives (e.g., solar optimization, pedestrian
movement) could enhance its applicability.

Incorporating additional environmental factors like thermal comfort,
pollutant dispersion, and energy consumption into the optimization
framework could enhance its applicability and provide more comprehensive
design guidelines.

\section{Conclusion}

This proof-of-concept study presents a novel, automated framework for
discovering climate-resilient urban design guidelines. By leveraging QD
optimization and RANS simulations, we demonstrate the potential to
uncover both conventional and unexpected urban planning strategies,
providing a foundation for further refinement and real-world
application.

\bibliographystyle{apalike}
\bibliography{bib}

\end{document}

%%